% RESEARCH_DRAFT_SHA256: a603f073cc377897730649d9efe3534813a4e525c2620f898e48de6b3505184e
% Layout and source-reference wrappers expanded for the web renderer.
\documentclass[11pt]{article}
\usepackage[margin=25mm]{geometry}
\usepackage{fontspec,unicode-math}
\setmainfont{Latin Modern Roman}
\setmathfont{Latin Modern Math}
\usepackage{amsmath,amssymb,mathtools,xurl,hyperref}
\hypersetup{colorlinks=true,urlcolor=blue,linkcolor=blue}
\setcounter{secnumdepth}{0}
\setlength{\parindent}{0pt}
\setlength{\parskip}{5pt}
\setlength{\emergencystretch}{3em}
\newcommand{\sref}[2]{\hyperlink{source:#1:#2}{[S#2]}}
\newcommand{\cataloguescope}[1]{\paragraph{Source record.}#1}
\newcommand{\floor}[1]{\left\lfloor#1\right\rfloor}
\begin{document}
\section*{Research-attempt provenance}
This standalone research attempt retains the complete original definition body below,
including its historical source and preparation notes. Those notes describe the
earlier catalogue preparation; the new attempt and its verification are identified
separately after that body. No original ID, catalogue status, or source statement is changed.

% UnsolvedMath ID 30003654; source catalogue code OWR-15957-007
\section{Certificates of Polynomial Nonnegativity in Copositive Optimization}\label{30003654}
{\small UnsolvedMath ID 30003654 · Computer Science}
\subsection{Definitions and mathematical statement}

\paragraph{Shared notation.}Standard number systems are $\mathbb N,\mathbb Z,\mathbb Q,\mathbb R,\mathbb C$. Specialized notation and index ranges follow the cited source and the definitions below.


The principal ideal $(h)$ consists of polynomial multiples of $h$. A nonnegativity certificate expresses a target polynomial through the specified nonnegative polynomial cone and polynomial multiples of equations defining the feasible variety.
\paragraph{Record scope and qualifications.}
The catalogue title is retained as metadata. Named certificates, contracts, bounds or frameworks in that title are not asserted as results of the original source extract.
The exact certificate identity and ambient nonnegative cone are not written in the supplied question.

Source-record limitation: The certificate identity and ambient nonnegative polynomial cone referred to by the question are absent.
\paragraph{Statement recorded in the supplied source.}
For polynomials nonnegative on an unbounded set $S\cap h^{-1}(0)$, under what hypotheses can certificates expressed through polynomials nonnegative on $S$ and the ideal $(h)$ be used to extend copositive-polynomial formulations and polynomial-optimization methods? \href{https://doi.org/10.4171/owr/2017/52}{[S2]}
\paragraph{Source status.}
The supplied catalogue labels this source record open; that metadata was checked against the supplied snapshot on 2026-10-11. The label is preserved for provenance and is not an independent verification of present open status or mathematical correctness. The catalogue literature assessment is dated 2026-08-21.
\subsection{Short English statement}
Characterize when the stated certificates support polynomial optimization on unbounded constrained sets.
\subsection{Sources}
\begin{itemize}
\item[S1] Ulam AI and the UnsolvedMath Contributors, UnsolvedMath: A Curated Collection of Open Mathematics Problems, record 30003654 (OWR-15957-007). Catalogue curation and metadata: CC BY 4.0. The source extract is retained; definitions, English summaries and qualifications are editorial additions.
\url{https://huggingface.co/datasets/ulamai/UnsolvedMath}.
\item[S2] Copositivity and Complete Positivity (2018), Oberwolfach Reports, DOI 10.4171/owr/2017/52, pp. 3071--3120
\url{https://doi.org/10.4171/owr/2017/52}.
\item Source-specific Oberwolfach Reports licences, including CC BY-SA 4.0 where applicable, are retained; see THIRD\_PARTY\_NOTICES.md.
\end{itemize}
\paragraph{Preparation and review.}
The original source and UnsolvedMath contributors retain attribution for the source extract. Codex (GPT-6; exact serving revision not exposed) assisted with definitions, English summaries, scope notes and formatting. No new human mathematical authorship or independently verified solution is asserted. Permanent source ID 30003654 and code OWR-15957-007 checked against the complete supplied registry and current repository. Distinct upstream records, including repeated questions and broad or damaged extracts, are retained for complete coverage of the source CS category. A repeated record is not a new independent conjecture. Complete mathematical-equivalence and current-literature review has not been performed. Preparation and final review follow contribution prompts C and D, with source gaps disclosed. Structural and rendering checks do not establish a new theorem or open status.
\label{end:30003654}

\section{Literature and context check}
Vera's joint-work talk with Zuluaga in OWR 52/2017, pp.~3108--3109, was read. It already states a certificate theorem under appropriate hypotheses and explicitly mentions closure in the compact case, then asks about consequences. The full theorem's hypotheses were not reconstructed from that short report. Root Codex GPT-6 supplied the auxiliary $h=x^2$ obstruction and the coordinate-hyperplane route; OpenAI gpt-6-astra independently checked them and developed the coefficient-growth observation below. No new general certificate theorem is claimed.
\section{Research attempt: Exact certificates, closure, and a projection-invariant case}

\paragraph{A precise exact-certificate test.}
Write $\operatorname{Pos}(S)$ for the real polynomials nonnegative at every point of a closed set $S$. An exact certificate of the type investigated here is
\[
 p=q+hr,\qquad q\in\operatorname{Pos}(S),\quad r\in\mathbb R[x_1,\ldots,x_n].
\]
Such a certificate implies $p\ge0$ on $S\cap\{h=0\}$ by substitution. Its unrestricted converse fails even on a very simple unbounded zero set. Take $S=\mathbb R^2$, $h(x,y)=x^2$ and $p(x,y)=x$. The polynomial $p$ vanishes on $\{x=0\}$. If $p=q+x^2r$ with $q\ge0$ globally, then for every fixed $y$, $q(0,y)=0$ is a minimum of the differentiable function $x\mapsto q(x,y)$, so $\partial_xq(0,y)=0$. But $q=x-x^2r$ gives $\partial_xq(0,y)=1$, a contradiction.

\paragraph{Closure repairs this example but loses bounded coefficients.}
For every $\delta>0$,
\[
 x+\delta=\frac{(x+2\delta)^2}{4\delta}-x^2\frac1{4\delta}.
\]
Thus $x$ is a coefficientwise limit of polynomials in $\operatorname{Pos}(\mathbb R^2)+(x^2)$, even using polynomials of degree at most two. This does not contradict a theorem formulated with closure. It also shows why compactness of certificate parameters cannot be inferred from a degree bound: the displayed multiplier coefficient diverges.

More generally, for any fixed degree bound on $r_\delta$, every exact family $x+\delta=q_\delta+x^2r_\delta$ with $q_\delta\ge0$ must have unbounded multiplier coefficients along $\delta\downarrow0$. Otherwise some sequence tending to zero would have a coefficientwise convergent subsequence of multipliers in the finite-dimensional bounded coefficient space. Then $q_\delta=x+\delta-x^2r_\delta$ would converge pointwise to a globally nonnegative polynomial $x-x^2r$, yielding the forbidden exact certificate. This proves the obstruction without assuming that the particular displayed family is optimal.

\paragraph{A sufficient hypothesis with an explicit formula.}
In contrast, let $h(x,y)=x$, where $y$ may denote any finite tuple of other variables, and suppose $S$ is invariant under the projection $(x,y)\mapsto(0,y)$. If $p\ge0$ on $S\cap\{x=0\}$, set $q(x,y)=p(0,y)$. The projection hypothesis implies $q\ge0$ on all of $S$. Expanding $p$ in powers of $x$ shows that $p-p(0,y)$ is divisible by $x$, so
\[
 p(x,y)=p(0,y)+x r(x,y)
\]
is an exact certificate. If $p$ has degree $d\ge1$, the constructed multiplier has degree at most $d-1$; if $p$ is independent of $x$, take $r=0$. This is a sufficient condition, not a necessary one or an algorithm for testing arbitrary polynomial nonnegativity on $S$.

The examples $h=x$ and $h=x^2$ have the same zero set but different principal ideals and different exact-certificate behavior. Therefore any proposed characterization based solely on the underlying zero set, ignoring the generator/ideal or the role of closure, cannot justify the exact identity under discussion. Passing from this restricted formula to effective copositive optimization hierarchies requires additional degree, positivity and convergence arguments.
\section{Verification and adversarial audit}
The derivative-at-a-minimum contradiction, closure identity, coefficient compactness argument and monomial-division formula were independently checked. The projection hypothesis is used exactly where nonnegativity is extended from the hyperplane to $S$. The root's mathematical contribution is expressly credited; no independent human review or full audit of the source theorem is asserted.
\paragraph{Audit corrections.} Separated exact identities from coefficientwise closure and distinguished the ideal $(x)$ from $(x^2)$ despite their identical zero sets. Added a proof that bounded-degree closure approximations need unbounded multiplier coefficients in the example.
\paragraph{External computation.} None; symbolic arguments and boundary checks only.
\section{Scope and remaining gaps}
The report's actual general hypotheses, closure topology, resulting optimization hierarchy and convergence guarantees remain untreated. The exact certificate is proved only for the stated projection-invariant hyperplane case, while the nonradical example rejects an unconditional converse. Neither refutes the source's qualified theorem nor resolves its broad question about extensions.
\paragraph{Integration review.} Codex (GPT-6; exact serving revision not exposed) reviewed the written arguments and expanded layout and source-reference wrappers. Repository metadata and the production web renderer were checked using Python 3.13 and JavaScript with MathJax 3.2.2. This model-assisted review is not independent human proof verification; no PDF compilation or formal proof checking is asserted.

\section{Completion Estimate}
\noindent\textbf{COMPLETION: 0\%}
The calculations below the preserved source statement investigate a concrete route,
but establish no substantiated new progress toward the entire original target.
Known special cases and failed sufficient conditions are not credited as a solution.
\section{References}
The original source references and catalogue attribution are preserved in the
Sources section above. The following URLs were actually consulted for this attempt;
retention of an additional original reference does not claim fresh consultation of it.
\begin{itemize}
\item \url{https://ems.press/content/serial-article-files/46715} (consulted 2026-10-11). 
\end{itemize}
\section{Model and contributor disclosure}
OpenAI gpt-6-astra, requested reasoning configuration ultra, generated this research
attempt on 2026-10-11. The exact serving revision is not exposed. Codex GPT-6
performed orchestration, repository-format review, and written-proof review of this attempt; its exact serving revision
is also not exposed. The mathematical route, boundary checks and final adversarial
self-audit followed contribution prompts A then B in
\texttt{docs/contribution\_prompts.md}, after reading \texttt{CONTRIBUTING.md}.
Python 3.14.2 assembled the standalone documents; any mathematical computation is
specified in Verification. No theorem prover was used. Codex GPT-6 also supplied the nonradical exact-certificate example and the projection-invariant hyperplane construction; gpt-6-astra independently checked and expanded these mathematical suggestions. No human mathematical
authorship, substantial human proof correction, or independent human proof checking
is asserted. Model-assisted review, source retrieval, and structural checks do not
establish correctness or novelty. Full private reasoning transcripts are unavailable
as reproducibility artifacts; the complete arguments, inputs and outcomes actually
used are stated here or in the companion batch audit.
\section{Final status}
Unresolved as a new research contribution. This is a concrete route investigation,
not a claim of new complete resolution or a current-status certification. A known
answer reproduced for context does not become a new contribution. The catalogue's
source-status label is preserved; this attempt status does not assert that a
source-labelled solved question is open.
\end{document}
